% C.3ao reconstructed after the 2026-11-19 scratch loss from retained proof text.
\paragraph{C.3ao (Refuted, with specified counterexamples): ordinary restriction injectivity in all remaining higher dimensions.}
For every $r\ge2$ and $N\ge2r+1$, we construct an excellent
normal strictly henselian local complex ring $R$ of dimension
$N$, with regular puncture $X$ and regular punctured principal
divisor $Y$, and a specified infinite-order class
\[
 \gamma\in\ker\!\left(
 \mathrm{CH}^r(X)_{\mathrm{alg}}
       \longrightarrow\mathrm{CH}^r(Y)_{\mathrm{alg}}\right).
\]
Algebraic equivalence is the ordinary relation using actual
families on the fixed product. All differentials in the
detecting calculation are relative to a stated subfield
$k_0\subset\mathbb C$, rather than to $\mathbb C$.

\begin{proof}
\emph{1. The ring and the smooth divisor.}
Put $m=N-r-1\ge r$. For $0\le j<r$, let
\[
 C_j=\{z_{j,0}^4+z_{j,1}^4+z_{j,2}^4=0\}
          \subset\mathbb P^2_{\overline{\mathbb Q}}.
\]
We use the canonical identification $K_{C_j}=\mathcal O_{C_j}(1)$
given by the residue convention. On $z_{j,2}=1$, with
$u=z_{j,0}$ and $v=z_{j,1}$, choose its basis
\[
 \omega_{j,0}=\frac{du}{4v^3},\qquad
 \omega_{j,1}=\frac{u\,du}{4v^3},\qquad
 \omega_{j,2}=\frac{v\,du}{4v^3}.
\]
These are the three canonical sections corresponding to
$z_{j,2},z_{j,0},z_{j,1}$. Let
\[
 S=\prod_{j=0}^{r-1}C_j,\qquad
 M=\boxtimes_{j=0}^{r-1}K_{C_j}.
\]
Choose the $(r+1)3^r$ numbers
\[
 \theta_{i,\mathbf k}\quad
 (0\le i\le r,\ \mathbf k\in\{0,1,2\}^{r})
\]
algebraically independent over $\overline{\mathbb Q}$.
Let $k_0$ be the algebraic closure of their generated
field inside $\mathbb C$, and set
\[
 s_i=\sum_{\mathbf k}\theta_{i,\mathbf k}
                 \bigotimes_{j=0}^{r-1}\omega_{j,k_j}
       \in H^0(S_{k_0},M),\qquad 0\le i\le r.
\]
These sections have no common zero. Indeed $M$ is generated
by its displayed $3^r$ sections. Over the affine space of
$(r+1)$ tuples of sections, the incidence that a point of
the $r$-dimensional projective variety $S$ is a common zero
has dimension at most one less than the parameter space:
evaluation imposes $r+1$ independent linear equations.
Its image is closed by properness and is therefore proper.
It is defined over $\overline{\mathbb Q}$, so the specified
independent coefficients avoid it.

After extension to $\mathbb C$, put
\[
 V=\mathbb P^m\times S,\qquad
 L=\mathcal O_{\mathbb P^m}(1)\boxtimes M,\qquad
 f=\sum_{i=0}^{r}x_i s_i\in H^0(V,L).
\]
The unused coordinates $x_{r+1},\ldots,x_m$ are allowed.
The surjection $\mathcal O_S^{m+1}\to M$ with entries
$s_0,\ldots,s_r,0,\ldots,0$ has a vector-bundle kernel
of rank $m$. Thus $H=V(f)\subset V$ is its projective
bundle of lines, smooth and connected. Let
\[
 B=\bigoplus_{n\ge0}H^0(V,L^n),\qquad
 R=(B_{B_+})^{\mathrm{sh}},\qquad
 X=\operatorname{Spec}R\setminus\{\mathfrak m_R\},\qquad
 Y=X\cap V(f).
\]
The residue field in the strict henselization is the
specified $\mathbb C$. The Segre section ring $B$ of the
projectively normal plane quartics and projective space
is a normal domain generated in degree one. Its punctured
cone is the complement of the zero section in
$\operatorname{Tot}(L^{-1})$, hence smooth. Therefore $R$
is excellent and normal of dimension $m+r+1=N$, and
$X,Y$ are regular. The nonzero element $f$ is a
nonzerodivisor. As in C.3i, the normalization of $B/(f)$
is the full section ring of $(H,L_H)$; its defect is
supported at the vertex and is finite-dimensional.
Indeed the section exact sequence and Serre vanishing
give precisely this assertion. Its graded vertex is
the unique point over the original vertex. Thus its
strict-local puncture is the same $Y$.

Let $Z\to\operatorname{Spec}R$ be the base change of
the blow-up of the cone vertex. It is regular, with
exceptional divisor $V$, projection $q:Z\to V$, and
\[
 I_V=\mathfrak m_R\mathcal O_Z=\mathcal O_Z(-V),\qquad
 N_{V/Z}=L^{-1},\qquad Z\setminus V=X.
\]
Its completion along $V$ is the split formal line bundle
with conormal $L$. In particular its first infinitesimal
neighborhood is
\[
 V_2=\operatorname{Spec}_V(\mathcal O_V\oplus L),\qquad L^2=0.
\]

\emph{2. An actual line bundle in the full Picard kernel.}
Choose $\xi_0\in H^1(C_0,\mathcal O_{C_0})$ Serre-dual
to $\omega_{0,0}$ in the displayed basis, and pull it
back to $V$. We specify
\[
                       a=f\xi_0\in H^1(V,L).
\]
If $\eta_0\in H^1(C_0,K_{C_0})$ has trace one, the
K\"unneth formula writes this as
\[
 a=\eta_0\otimes
   \sum_{i=0}^{r}x_i
    \sum_{\mathbf k'\in\{0,1,2\}^{r-1}}
       \theta_{i,(0,\mathbf k')}
           \bigotimes_{j=1}^{r-1}\omega_{j,k_j}.
\]
For $n\ge2$, $H^1(V,L^n)=0$: Serre duality gives
$H^1(C_j,K_{C_j}^n)=0$, and positive powers on
projective space have no higher cohomology.

The characteristic-zero logarithm gives a formal line
bundle with trivial restriction to $V$, logarithmic
class $a$ in degree one, and zero higher graded classes.
Explicitly, a \v Cech cocycle for $a$ gives formal
transition functions $\exp(a_{\alpha\beta})$; every
finite level uses a finite exponential polynomial.
Split units and surjective transition maps on global
nilpotent units give
\[
 \operatorname{Pic}(\widehat Z)
       =\operatorname{Pic}(V)\oplus H^1(V,L).
\]
This is the formal calculation of C.3v with the stated
K\"unneth vanishing.

Proper formal existence algebraizes this line bundle
on $Z_{\widehat R}$. Elkik's punctured henselian
comparison supplies its unique class
$\mathcal A_X\in\operatorname{Pic}(X)$
\citep[Theorem 30.27.1, Tag 088E]{rootStacksExistence}
\citep[Section III, Theorem 3 and its Picard corollary,
pp.~569,573]{D_Elkik1973}.
Extend a Cartier divisor representing $\mathcal A_X$
across $V$ on the regular scheme $Z$. Its completed
extension differs from the chosen one by an integral
exceptional divisor. An exceptional twist therefore
gives an actual line bundle $\mathcal A$ on $Z$
with exactly the formal class $(0,a)$. In particular
\[
                         \mathcal A|_V\simeq\mathcal O_V.
\]
The exact sequence
\[
 0\longrightarrow\mathcal O_V\xrightarrow{\ f\ }L
       \longrightarrow L_H\longrightarrow0
\]
shows that $a$ restricts to zero in $H^1(H,L_H)$.
Restriction of the split formal bundle respects the
grading and logarithm. A degree-one \v Cech coboundary
exponentiates to a trivialization at every order;
possible higher $H^1(H,L_H^n)$ contribute no hidden
component. Thus the completed restriction to $Y$ is
trivial. The injective Elkik comparison for the normal
section-ring model of the same puncture gives
\[
                           \mathcal A_X|_Y\simeq\mathcal O_Y
                           \quad\hbox{integrally}.
\]
These are actual line bundles, not merely tangent
vectors of local Picard functors.

\emph{3. An ordinary cycle and its absolute cycle map.}
Choose $u_1,\ldots,u_{r-1}\in\mathbb C$ algebraically
independent over $k_0$, choose $v_j^4=-1-u_j^4$, and put
\[
 p_j=(u_j:v_j:1),\qquad
 o_j=(1:\epsilon:0),\quad\epsilon^4=-1,\qquad
 \lambda_j=\mathcal O_{C_j}(p_j-o_j).
\]
Use a fixed algebraic root $\epsilon$. Let $q_j:Z\to C_j$
be the projections and define
\[
 \widetilde\gamma
      =c_1(\mathcal A)\prod_{j=1}^{r-1}c_1(q_j^*\lambda_j)
             \in\mathrm{CH}^r(Z),\qquad
 \gamma=\widetilde\gamma|_X.
\]
Varying $p_1$ on the fixed proper curve $C_1$ by its
diagonal divisor gives an actual ordinary family on
$X\times C_1$; its difference at $p_1,o_1$ is $\gamma$.
Thus $\gamma$ is ordinarily algebraically trivial.
The trivial restriction of $\mathcal A_X$ gives
$\gamma|_Y=0$ integrally.

For a smooth finite-type complex scheme $U$, the integral
Milnor Bloch formula and the symbol map give
\[
 \mathrm{CH}^r(U)=H^r(U,\mathcal K^M_r)
       \longrightarrow H^r(U,\Omega^r_{U/k_0}),\qquad
 \{g_1,\ldots,g_r\}\longmapsto
                   d\log g_1\wedge\cdots\wedge d\log g_r.
\]
The Steinberg relation maps to zero since
$d\log g\wedge d\log(1-g)=0$. The integral formula is
\citet[Theorems 7.1 and 7.5, author pp.~20--22]{cyKerz2009Milnor};
compatibility with divisor intersections is the cup and
concatenation formula of
\citet[Section 3, pp.~557--558]{cyBraunling2013}.
The image of a divisor product is therefore the product
of its $d\log$ first Chern classes. Those classes and
their \v Cech products are functorial for arbitrary
pullbacks, including a nonreduced target
\citep[Section 50.9 and Lemma 50.9.1, Tag 0FMD]{cyStacksHodgeChern}.

Any proposed rational equivalence for the displayed
class on $Z$ descends to a finite pointed etale
neighborhood of the cone. Its pulled-back resolution
is smooth of finite type over $\mathbb C$. We can thus
pull the preceding class to $V_2$ and then tensor its
differential sheaf with $\mathcal O_V$. This detects
rational nonvanishing on the strict-local $Z$ itself;
we do not assert a Bloch formula on $V_2$.
The split square-zero algebra gives
\[
 \Omega^1_{V_2/k_0}\otimes\mathcal O_V
       =\Omega^1_{V/k_0}\oplus L,
 \qquad
 \Omega^r_{V_2/k_0}\otimes\mathcal O_V
       =\Omega^r_{V/k_0}\oplus
                         L\otimes\Omega^{r-1}_{V/k_0}.
\]
In local notation $t^2=0$, the transition $1+t a_{\alpha\beta}$
has $d\log$ class $a_{\alpha\beta}dt$ after tensoring
with $\mathcal O_V$. Thus $c_1(\mathcal A)$ has
conormal component $a$. Notice that tensoring the
differential sheaf retains $dt$; ordinary differential
pullback to $V$ would instead kill that component.
The product consequently has conormal class
\[
                   \pm a\smile\prod_{j=1}^{r-1}
                                  c_1^{\mathrm{abs}}(\lambda_j).
\]
Because $V$ and all its displayed line bundles are defined
over $k_0$, the splitting
\[
 \Omega^1_{V/k_0}
   =\Omega^1_{V/\mathbb C}
          \oplus\mathcal O_V\otimes_{\mathbb C}
                         \Omega^1_{\mathbb C/k_0}
\]
is canonical. Project the last class to
\[
             H^r(V,L)\otimes_{\mathbb C}
                             \Omega^{r-1}_{\mathbb C/k_0}.
\]

\emph{4. The nonzero residue calculation.}
For a curve defined over $k_0$, a point $p$ with local
coordinate $u(p)=b$, and an algebraic base point $o$,
the line bundle $\mathcal O(p-o)$ has local gluing
function $u-b$. Its constant differential component is
$-db/(u-b)$. Under Serre duality, pairing with a regular
form $\omega=g(u)du$ gives the residue
\[
          -g(p)\,db=-p^*\omega
                        \quad\hbox{in }\Omega^1_{\mathbb C/k_0}.
\]
Here $p:\operatorname{Spec}\mathbb C\to C_{k_0}$,
so the notation on the right includes the field
differential. The contribution of $o$ is zero, and
the relative complex first Chern class is zero because
the divisor has degree zero.

For our residue basis, the resulting conormal component
is, up to a single nonzero sign,
\[
 (\eta_0\otimes\cdots\otimes\eta_{r-1})\otimes
 \left(\sum_{i=0}^{r}x_i T_i(p)\right)
 \frac{du_1\wedge\cdots\wedge du_{r-1}}
                   {4^{r-1}\prod_{j=1}^{r-1}v_j^3},
\]
where every $\eta_j\in H^1(C_j,K_{C_j})$ has trace one and
\[
 T_i(p)=\sum_{\mathbf k'}\theta_{i,(0,\mathbf k')}
                  \prod_{j=1}^{r-1}b_{j,k_j},
       \qquad (b_{j,0},b_{j,1},b_{j,2})=(1,u_j,v_j).
\]
The canonical tensor section $T_0$ is nonzero, since
the displayed tensors form a basis and its coefficients
are not all zero. The specified point of the product
of curves is geometric generic over $k_0$, so $T_0(p)\ne0$.
The forms $du_1,\ldots,du_{r-1}$ are independent over
$\mathbb C$: the $u_j$ extend to a transcendence basis
of $\mathbb C/k_0$, and characteristic zero makes the
remaining algebraic extension separable. Also $v_j\ne0$.
The displayed class is therefore nonzero. It follows
that $\widetilde\gamma$ has infinite order in
$\mathrm{CH}^r(Z)$.

\emph{5. Excluding the exceptional boundary.}
Suppose $\gamma=0$ in $\mathrm{CH}^r(X)_{\mathbb Q}$.
Chow localization gives
\[
               \widetilde\gamma=i_*\beta,
               \qquad\beta\in\mathrm{CH}^{r-1}(V)_{\mathbb Q}.
\]
Since $\mathcal A|_V$ is trivial, refined pullback and
self-intersection give
\[
                  0=i^*\widetilde\gamma=-c_1(L)\beta.
\]
Put $h=c_1(\mathcal O_{\mathbb P^m}(1))$ and
$b=c_1(M)$. By the projective-bundle formula write
\[
                  \beta=\sum_{j=0}^{r-1}h^j\beta_j,
              \qquad\beta_j\in\mathrm{CH}^{r-1-j}(S)_{\mathbb Q}.
\]
Because $m\ge r$, multiplication by $h+b$ is not
truncated in any of the powers $h^0,\ldots,h^r$.
The coefficient of $h^r$ first gives $\beta_{r-1}=0$.
Descending through the coefficients gives successively
$\beta_{r-2}=\cdots=\beta_0=0$. This recursion even
proves the corresponding injectivity integrally; it
uses no conjecture about Chow Lefschetz maps.
Thus $\widetilde\gamma=0$ rationally on $Z$, contrary
to the absolute differential detector. Hence $\gamma$
has infinite order on $X$ and is the required ordinary
restriction kernel class.
\end{proof}

The finite certificate checks the defining equations,
dimension and incidence counts, projective-coefficient
recursion, canonical-tensor normal forms and the exact
local residue and conormal identities. The generic
nonvanishing, actual algebraization and Chow argument
are proved above, not inferred from those finite checks.
